\Needspace{8\baselineskip}
\section{Concrete integration into Forge, Leant, and Djex}
\label{sec:integration}

\subsection{Reuse the existing controller; add three semantic front ends}

Forge already has proposed obligation, candidate, and conditional-candidate
contracts. In particular, its runtime specimen has a
\code{ConditionalCandidate} carrying premises and a constructor term, and an
\code{EngineReply} that can return progress or a reasoned failure~\cite{forge-runtime}.
The workers in this report fit that contract. They do not require a replacement
for Forge's scheduler or an embedded second Lean kernel.

The missing layer is a \emph{proved extraction} from a source expression into
one of the three fragments. A certificate about an independently invented
transition system is not a proof of the user's recursive function. The
following proposed module split isolates this work:

\par\noindent\begin{minipage}{\linewidth}
\begin{lstlisting}
Forge/Behavior/PolynomialState.lean   -- source transition descriptions
Forge/Behavior/Weighted.lean          -- affine/linear fold descriptions
Forge/Behavior/BinomialSum.lean       -- summand and finite-support bridges
Forge/Certificate/InductiveIdeal.lean -- finite identities and soundness
Forge/Certificate/Weighted.lean       -- closure matrices and word replay
Forge/Certificate/Telescoping.lean    -- regularized summation, uniqueness
Forge/Worker/InductiveIdeal.lean      -- problem extraction and CAS adapter
Forge/Worker/Weighted.lean            -- exact reachable-space search
Forge/Worker/Telescoping.lean         -- coefficient search and proof assembly
Forge/Export/BehaviorProof.lean       -- stable theorem/lemma output
\end{lstlisting}
\end{minipage}\par

These are proposed files, not files asserted to exist in the reviewed
repository. The supplied Python files implement the corresponding search and
finite replay functions; the Lean source front ends remain to be implemented.

\subsection{Polynomial extraction and replay}

For the first version, accept explicit polynomial-state descriptions rather
than arbitrary recursive programs. A description contains an exact coordinate
telescope, initial polynomial maps, update polynomials, equality-guard
polynomials, a goal polynomial, and Lean proof obligations tying these to the
source.

For an ordinary recursive function, derive its transition view from checked
equations. A natural-number loop can package its loop index into the state,
for example $(i,s)\mapsto(i+1,F(i,s))$. A list fold can use one update per
finite letter class. The equation compiler must use simultaneous substitution:
under $(x,y)\mapsto(y,x)$, both right sides refer to the old state. A sequential
substitution routine would silently implement a different transition.

Mathlib's \code{MvPolynomial} evaluation interface provides
\code{eval}, \code{eval₂}, \code{eval₂Hom}, and \code{aeval}; evaluation and
substitution lemmas support the mathematical interpretation of coefficient
identities~\cite{mathlib-mvpoly}. A replay theorem can evaluate each serialized
polynomial in the source coordinates, use the checked identities to prove the
invariant conjunction's step, and apply ordinary reachability induction.
The core-only \code{reachableInvariant} specimen in this package shows the
last induction principle independently of any polynomial implementation.

There are two practical replay implementations. For a narrow first release,
translate each coefficient list to an ordinary Lean polynomial expression,
prove the finite identities with \code{ring} and rational normalization, and
apply the generic soundness theorem. For a faster later release, implement a
small reflected sparse-polynomial evaluator and prove its correctness once.
Both routes end in a Lean proof term. The latter is an optimization of repeated
identity checking, not a prerequisite for testing the mathematical fragment.

Avoid making a proof of Gr\"obner-basis correctness an early dependency. The
final certificate never requires that fact. It only requires original-generator
identities. The uncompiled \code{ReplayTargets.lean} file includes the concrete
nonlinear-curve identity and its invariant step, illustrating the small output
that the search can leave behind.

\subsection{Weighted extraction and replay}

Start with an explicit weighted representation over \code{Rat}. Its dimensions
should be type-indexed by \code{Fin N} and \code{Fin r}, not represented by
unchecked array lengths inside the mathematical theorem. The external decoder
first establishes the lengths, then constructs the corresponding Lean values.
The theorem consumes $B,c_0,C_a$ and proofs of the three identities in
Equation~\eqref{eq:weighted-cert}.

Mathlib documents \code{Matrix.vecMul}, \code{Matrix.mulVec}, and associativity
lemmas including \code{Matrix.vecMul_vecMul}~\cite{mathlib-matrix}. These give a
straightforward row-vector proof. The supplied \code{simulateWord} specimen
abstracts the argument: if an embedding commutes with each letter step, then it
commutes with running any word. For the matrix certificate, that embedding is
$c\mapsto cB$, and the commuting square is $BM_a=C_aB$.

The first source front end should support a finite enumerated alphabet and
explicit tuples of rational affine accumulators. Syntactic recognition can
propose the matrices, but Lean must prove the source step equals the proposed
affine map for every letter and state. After that bridge is established, the
worker can compare independently structured folds, not just normalized copies
of one expression.

The negative route returns a concrete word and a proved inequality between its
source observations. For source programs with executable data, this can be
shown as a diagnostic even when the goal itself is not a proposition asserting
universal equality. It does not authorize a global non-inhabitation conclusion
about an unrelated synthesis type.

\subsection{Summation extraction and replay}

The initial sum front end should recognize exactly
\[
 \sum_{k\in\operatorname{range}(n+1)} (\operatorname{Nat.choose}(n,k):\Q)^m
\]
for a literal supported $m$. It can later add source-equivalent spellings by
proved rewrites. Recognition must preserve the cast location: a natural power
cast into $\Q$ can be related to the rational power by a theorem, but neither
expression should be guessed from surface syntax.

Useful Mathlib ingredients include \code{Nat.choose_eq_zero_of_lt},
\code{Nat.choose_pos}, \code{Nat.choose_symm}, and
\code{Nat.choose_mul_factorial_mul_factorial}~\cite{mathlib-choose}. The last is
especially useful for cross-multiplied ratio proofs, avoiding premature field
division. Finite-sum lemmas in the \code{Finset} big-operator API provide the
reindexing and induction infrastructure~\cite{mathlib-finset}. The package does
not assume an existing theorem with exactly the pole-safe certificate signature;
that theorem is one of the explicit implementation tasks.

Prove the regularized summation principle once for arbitrary $r,m,p,U$ satisfying
the finite obligations. Then a particular certificate supplies only rational
polynomial identities, fixed boundary-index calculations, positivity data, and
initial values. To prove equality with another expression, attach a second
recurrence proof and invoke the shared uniqueness principle. Do not recompute
the entire sum for a large numerical range during Lean replay.

The concrete Lean specimen includes the first-power, square, and cube
\emph{interior polynomial identities}. It does not contain a full Lean proof of
the binomial support bridges or Equation~\eqref{eq:franel}. Keeping this
boundary explicit prevents the short algebraic specimen from being mistaken
for a completed formalization of the sum theorem.

\subsection{A practical worker transaction}

A complete worker invocation should behave as follows. The frontend first
constructs a typed fragment problem and the bridge obligations. It calls a
search backend with the problem's immutable coefficient data. The result is
decoded into finite certificate data. The appropriate replay theorem proves
the fragment result. Finally, a proof constructor combines that result with
\emph{all} bridge proofs to produce the original goal.

\par\noindent\begin{minipage}{\linewidth}
\begin{lstlisting}
runBehaviorWorker(sourceGoal):
    fragment, bridges := extractSupportedFragment(sourceGoal)
    candidate := search(fragment)
    if no candidate: return NoResult(reason)
    finiteProof := replayFiniteCertificate(fragment, candidate)
    if bridges are already proved:
        return Candidate(transport(finiteProof, bridges))
    return ConditionalCandidate(
        premises = bridges,
        constructor = transportConstructor(finiteProof))
\end{lstlisting}
\end{minipage}\par

This is pseudocode for the proposed adapter, not a claim about an existing API
function named \code{runBehaviorWorker}. Forge's current controller vocabulary
already has the needed conditional shape~\cite{forge-runtime}. The new content
is the fragment extraction and the replay theorems, not the existence of a
conditional work item.

The result should usually be exported as a local or named lemma that ordinary
\grind{} can consume. The current Lean reference describes \grind{} as a
combination of equality and theory reasoning producing proof terms~\cite{lean-grind}.
The new lemma changes the facts available to that engine: a global invariant,
a fold equivalence, or a sum recurrence is information that a larger local
arithmetic budget would not necessarily discover.

\subsection{Concrete library choices}

\begin{longtable}{@{}p{0.20\textwidth}p{0.34\textwidth}p{0.40\textwidth}@{}}
\toprule
Layer & Choice & Exact role and boundary \\
\midrule\endhead
Prototype search & Python 3.13.5; SymPy 1.14.0 & Polynomial expansion, rational
coefficient matrices, exact nullspaces; tracked Buchberger search supplied in
this package. Executed. \\
Independent replay & Python standard library \code{fractions.Fraction} &
Canonical rational and sparse polynomial data, matrix arithmetic, exact finite
checks. Runs with \code{python -S}. Executed. \\
Larger ideal search & Sage/Singular & Use \code{I.groebner_basis()} and
\code{f.lift(I.gens())}; export original-generator multipliers and verify their
identity. Proposed, not executed here~\cite{sage-ideals}. \\
Larger linear algebra & FLINT-backed exact arithmetic & Replace repeated
rational solves with fraction-free or modular computations, then reconstruct
and recheck the same closure identities. Proposed~\cite{flint}. \\
Broader summation & SymPy Gosper; Sage \code{ore_algebra} & Antidifferences or
operator/certificate proposals; preserve source support and boundary proofs.
Proposed~\cite{sympy-concrete,ore}. \\
Lean replay & Lean and pinned Mathlib & Prove polynomial interpretation,
finite-space simulation, binomial support, telescoping, and uniqueness. Not run
in this package. \\
\bottomrule
\end{longtable}

A useful Sage adapter can be very small: create the exact rational polynomial
ring with a fixed variable order, construct $I$ from the original generators,
request \code{f.lift(I.gens())}, and serialize the returned polynomial list.
The documentation explicitly illustrates reconstructing $f$ by multiplying
that list by the generator list~\cite{sage-ideals}. Failure to lift is a search
failure; a malformed or incorrect returned identity is replay rejection.
There is no need to transmit a complete Gr\"obner trace to Lean.

The archive does not bundle these external systems or their licenses. Its own
code is released under MIT-0. Backend packaging and license compatibility
should be handled separately from the mathematical certificate format; the
format itself does not depend on the chosen search library.

\subsection{Leant and Djex: add behavior domains, not another candidate authority}

Djex's documented semantic layer consumes a source-owned typed candidate graph
rather than manufacturing a new term identity. Its current length domain
tracks a supported list-spine behavior through checked problem descriptions and
replay~\cite{djex-semantic}. Two additional domains fit this arrangement:

\emph{Polynomial-state contracts} assign a proved polynomial transition summary
to a source fold or recursive accumulator. The ideal worker can prove
postconditions about that state. \emph{Weighted-observation contracts} assign a
proved weighted representation to a scalar observable. The weighted worker can
prove equality or return a finite distinguishing input for that observable.

Each summary must carry its full source type, the exact chosen instances and
providers, and the theorem connecting the candidate's semantics to the
representation. These are not extra versions of the type certificate. They are
additional theorems about the same candidate. When a provider is polymorphic,
instantiate the summary at the same source-level type and universe selections
as the retained candidate graph; substituting a conveniently similar monomorphic
provider changes the problem.

The benefit is more than post-hoc scoring. A verified summary can discharge a
behavioral obligation for all inputs in its fragment. A counterexample can
reject a candidate quickly. But neither domain should prune unrelated
higher-order alternatives merely because their current partial terms lack an
extractable summary. Unsupported analysis remains absence of information.

\subsection{A targeted composition experiment for the remaining indexing gap}

Leant's current diagnosis is unusually concrete. The isolated integer-indexing
fold step is found at candidate 5, and the same step in a larger outer context
at candidate 219, while the original full indexing query still fails all 256
checked candidates in its recorded window~\cite{leant-indexing}. These are not
equivalent benchmark tasks, and the diagnosis explicitly leaves full indexing
unaccepted.

A focused follow-on experiment is to introduce a \emph{typed composition cut}.
When a candidate skeleton contains a function-valued fold, expose its step and
base as named synthesis obligations using their exact source types. Search the
step obligation first under a dependency-closed subset of the local telescope,
and retain a successful step as a candidate component. Reinsert it by applying
the original fold provider and independently checking the complete term.

The dependency-closed telescope is computable: start from variables occurring
in the hole's type, explicit behavioral contract, and selected provider
arguments; repeatedly add free variables occurring in the types and local
values of those variables, including instance dictionaries. Search this reduced
context as an additional lane, not as a completeness-preserving replacement for
the full-context lane. A solution found in the reduced context is valid in the
full context after checked weakening; failure in the reduced context proves
nothing about the larger query.

For the documented step type
\par\noindent\begin{minipage}{\linewidth}
\begin{lstlisting}
forall a. a -> a -> (Int -> a) -> Int -> a
\end{lstlisting}
\end{minipage}\par
a reusable component should still distinguish negative, zero, and positive
indices and retain the supplied default behavior. Cache the generalized checked
component under its complete source type and provider inventory. Do not reuse
an ordinal such as ``candidate 5'' as an identity, and do not assume the same
component cost in the outer search: usage history and pending holes can change
that cost.

This experiment is proposed, not implemented or accepted here. Its success
criterion is the \emph{original} indexing query at its recorded signature,
behavior observations, candidate window, and limits, with the exact displayed
term replayed. It is included because the repository now identifies a specific
composition bottleneck to investigate, rather than because another generic
suggestion to improve ranking would add much value.
